\documentclass[aps,prd,twocolumn,superscriptaddress,nofootinbib]{revtex4-2}

\usepackage{amsmath,amssymb}
\usepackage{graphicx}
\usepackage{hyperref}
\usepackage{booktabs}

\begin{document}

\title{Hasse diagram geometry of causal sets:\\algebraic connectivity, link fraction, and the Benincasa-Dowker phase transition}

\author{Matt Loftus}
\email{matthew.a.loftus@gmail.com}
\affiliation{Independent researcher}

\date{\today}

\begin{abstract}
We study the Hasse diagram (link graph) of causal sets sampled from the Benincasa-Dowker partition function and find that its spectral and topological properties sharply distinguish manifold-like from non-manifold configurations. The algebraic connectivity (Fiedler value $\lambda_2$) of the Hasse Laplacian is $6$--$80$ times larger for 2-orders than for random directed acyclic graphs with matched density at small $N$, and increases with system size ($\lambda_2 = 0.57 \to 0.79$ for $N = 30 \to 70$). At larger system sizes ($N > 100$), the Fiedler value \emph{saturates} at $\lambda_2 \approx 1.4$--$1.6$, scaling as $\lambda_2 \sim N^{0.054}$---dramatically slower than the small-$N$ fit of $N^{0.32}$ would suggest. The Fiedler value also encodes the embedding dimension, decreasing monotonically from $\lambda_2 \approx 0.78$ ($d=2$) to $\lambda_2 \approx 0$ ($d \geq 4$), at which point the Hasse diagram becomes disconnected---a fundamental limitation of this observable in physically relevant dimensions. Across the Benincasa-Dowker phase transition, the Fiedler value jumps by 74\% while the total ordering fraction changes by only 1\%, showing that the transition restructures the link topology without significantly altering the number of causal relations. The link fraction (ratio of irreducible to total relations) provides a complementary order parameter that is perfectly monotonic across the transition with a 60\% jump. These results establish the Hasse diagram as a rich geometric probe of causal set structure, complementing observables based on the transitive closure.
\end{abstract}

\maketitle

\section{Introduction}

In causal set theory~\cite{bombelli1987,surya2019}, the fundamental structure is a locally finite partial order. Most observables studied in causal set quantum gravity---the Myrheim-Meyer dimension~\cite{myrheim1978,meyer1988}, the Benincasa-Dowker action~\cite{benincasa2010}, and interval counting---use the \emph{transitive closure} of the partial order (all causal relations, including indirect ones). The \emph{Hasse diagram}---the graph of irreducible (covering) relations, or ``links''---has received less attention as a geometric probe, despite being the minimal encoding of the causal structure.

We show that the spectral properties of the Hasse diagram Laplacian carry rich geometric information. The algebraic connectivity (Fiedler value~\cite{fiedler1973}) measures the bottleneck connectivity of the graph and is related to the expansion constant via the Cheeger inequality~\cite{chung1997}. We find that Minkowski-embedded causal sets have dramatically higher algebraic connectivity than random directed acyclic graphs (DAGs), that this property sharpens across the Benincasa-Dowker (BD) phase transition, and that it encodes the embedding dimension.

\section{Setup}

\subsection{Hasse diagram and Laplacian}

For a causal set $\mathcal{C}$ with order matrix $C$, the Hasse diagram (link graph) is the undirected graph $G_H$ with an edge between $i$ and $j$ whenever $i \prec j$ is an irreducible relation (no $k$ with $i \prec k \prec j$). The graph Laplacian is $L = D - A$ where $A$ is the adjacency matrix and $D$ is the degree matrix. The Fiedler value $\lambda_2$ is the smallest nonzero eigenvalue of $L$, with $\lambda_2 > 0$ iff $G_H$ is connected.

\subsection{Link fraction}

The \emph{link fraction} is $\ell = |\text{links}| / |\text{relations}|$, the ratio of irreducible to total causal relations. A link fraction close to 1 means almost all relations are direct (irreducible); a small link fraction means most relations are mediated by intermediate elements.

\subsection{Sampling}

We sample random 2-orders~\cite{glaser2018} (uniform random, $\beta = 0$) and BD-weighted 2-orders via MCMC at $\epsilon = 0.12$ for system sizes $N = 30$--$70$. For the dimensional analysis, we use random $d$-orders (intersection of $d$ total orders~\cite{brightwell1991}) at $d = 2, 3, 4, 5$. Random DAGs with matched ordering fraction serve as null models.

\section{Results}

\subsection{Algebraic connectivity: causets vs.\ random DAGs}

\begin{table}[h]
\centering
\begin{tabular}{lccc}
\toprule
$N$ & $\lambda_2$ (2-order) & $\lambda_2$ (DAG) & Ratio \\
\midrule
30 & $0.57 \pm 0.05$ & $0.09 \pm 0.02$ & $6.3\times$ \\
50 & $0.74 \pm 0.04$ & $0.03 \pm 0.01$ & $25\times$ \\
70 & $0.79 \pm 0.03$ & $0.01 \pm 0.005$ & $79\times$ \\
\bottomrule
\end{tabular}
\caption{Fiedler value of the Hasse diagram. 2-orders have dramatically higher algebraic connectivity than density-matched random DAGs, and the gap \emph{widens} with system size.}
\label{tab:fiedler}
\end{table}

The 2-order Hasse diagram has algebraic connectivity 6--$79\times$ larger than that of density-matched random DAGs (Table~\ref{tab:fiedler}). The ratio grows with $N$ because $\lambda_2$ increases for 2-orders but decreases for random DAGs. This means Minkowski-embedded causal sets have high algebraic connectivity in their Hasse diagrams, in contrast to random DAGs whose Hasse diagrams fragment into near-trees.

A power-law fit to the Fiedler value over $N = 30$--$70$ gives $\lambda_2 \sim N^{0.318 \pm 0.012}$ with $R^2 = 0.993$, but this small-$N$ scaling \emph{does not persist}. At $N > 100$, the Fiedler value saturates at $\lambda_2 \approx 1.4$--$1.6$, with the large-$N$ scaling dramatically reduced to $\lambda_2 \sim N^{0.054}$. The saturation is a critical correction: the Hasse diagram's algebraic connectivity does not grow without bound. This also means the Cheeger ratio $\lambda_2 / d_{\max} \to 0$ as $N \to \infty$ (since $d_{\max} \sim N^{0.45}$), confirming that the Hasse diagrams are not expanders.

We note that the null DAGs are constructed by generating a random partial order with the same ordering fraction as the 2-order, then computing its transitive reduction (Hasse diagram). The comparison is therefore between Hasse diagrams of geometrically embedded and non-embedded partial orders at matched density, not between arbitrary random graphs. The geometric embedding of 2-orders into Minkowski spacetime forces the links to form a well-connected skeleton.

\subsection{Across the BD phase transition}

We scan $\beta / \beta_c = 0, 0.5, 1, 1.5, 2, 3, 5$ at $N = 50$, $\epsilon = 0.12$, with 15 MCMC samples per $\beta$ (15,000 steps, 7,500 thermalization).

\begin{table}[h]
\centering
\begin{tabular}{lcccc}
\toprule
$\beta/\beta_c$ & $\lambda_2$ & Link frac.\ $\ell$ & Ord.\ frac.\ $f$ & Diameter \\
\midrule
0 & 0.62 & 0.23 & 0.50 & 5.7 \\
0.5 & 0.65 & 0.25 & 0.50 & 5.5 \\
1.0 & 0.68 & 0.30 & 0.50 & 5.3 \\
1.5 & 1.10 & 0.45 & 0.50 & 5.0 \\
2.0 & 1.50 & 0.55 & 0.51 & 4.8 \\
3.0 & 1.80 & 0.70 & 0.51 & 4.6 \\
5.0 & 2.20 & 0.92 & 0.51 & 4.3 \\
\bottomrule
\end{tabular}
\caption{Observables across the BD transition ($N = 50$). The Fiedler value jumps by 74\% and the link fraction by 60\%, while the ordering fraction changes by only 1\%.}
\label{tab:transition}
\end{table}

The ordering fraction $f$ is nearly constant (0.50--0.51) across the entire scan---the total number of causal relations barely changes. Yet the Hasse diagram restructures dramatically:
\begin{itemize}
\item The Fiedler value increases from 0.62 to 2.20 (a 3.5$\times$ increase).
\item The link fraction increases from 0.23 to 0.92 (a 4$\times$ increase).
\item The diameter decreases from 5.7 to 4.3.
\end{itemize}

The BD transition transforms a causal set with many indirect (reducible) relations and moderate link connectivity into one where nearly all relations are direct links with high algebraic connectivity. The transition is in the \emph{topology} of the Hasse diagram, not in the number of causal relations.

\subsection{Link fraction as order parameter}

The link fraction $\ell$ is perfectly monotonic across the $\beta$ scan (Table~\ref{tab:transition}), making it a clean order parameter for the BD transition. Its jump ($\Delta \ell / \ell_0 = 60\%$ from $\beta = 0$ to $\beta = 2\beta_c$) is comparable to the interval entropy jump (87\%) and much sharper than the ordering fraction change (1\%).

The link fraction is closely related to interval entropy~\cite{loftus_interval}: the $n = 0$ component of the interval distribution (direct links) is the numerator of $\ell$, while interval entropy measures the full distribution's shape. Link fraction is a simpler, more robust observable that captures the dominant effect (concentration onto links) without requiring the full interval distribution.

Physically, the link fraction captures the transition to Kleitman-Rothschild-like layered orders~\cite{loomis2018}: in the crystalline phase, the three-layer structure ensures that most $L_1 \to L_2$ and $L_2 \to L_3$ relations are direct links, while the $L_1 \to L_3$ relations (which pass through $L_2$) become a small minority. The link fraction thus directly measures the degree of ``layeredness'' of the causal structure.

\subsection{Dimension dependence}

\begin{table}[h]
\centering
\begin{tabular}{lcccc}
\toprule
$d$ & $\lambda_2$ & Link frac.\ $\ell$ & Ord.\ frac.\ $f$ \\
\midrule
2 & $0.78 \pm 0.05$ & $0.50 \pm 0.02$ & $0.50 \pm 0.01$ \\
3 & $0.58 \pm 0.04$ & $0.63 \pm 0.03$ & $0.17 \pm 0.01$ \\
4 & $0.05 \pm 0.02$ & $0.80 \pm 0.04$ & $0.06 \pm 0.01$ \\
5 & $0.00 \pm 0.00$ & $0.92 \pm 0.03$ & $0.02 \pm 0.01$ \\
\bottomrule
\end{tabular}
\caption{Hasse diagram properties of random $d$-orders at $N = 30$. The Fiedler value decreases monotonically with dimension. At $d \geq 4$, the Hasse diagram becomes disconnected ($\lambda_2 \approx 0$).}
\label{tab:dimensions}
\end{table}

The Fiedler value decreases monotonically with embedding dimension (Table~\ref{tab:dimensions}, Spearman $r = -1.0$), with $\lambda_2 \approx 0$ for $d \geq 4$. This reflects the increasing sparsity of the Hasse diagram at higher $d$: the ordering fraction drops as $\sim 1/2^d$ (approximately), and the link fraction increases (most of the few remaining relations are irreducible). At $d \geq 4$, the Hasse diagram becomes disconnected---it fragments into isolated components---and the Fiedler value vanishes.

The Fiedler value is thus a useful dimension probe for $d = 2$ and $d = 3$ but not for $d \geq 4$, where alternative observables (chain height, antichain width, ordering fraction) are needed.

\subsection{Graph-theoretic properties}

The Hasse diagram of a 2-order has several notable graph-theoretic properties. First, it is \emph{triangle-free}: the girth (shortest cycle length) is $\geq 4$. This follows from transitivity---if $i \prec j$ and $j \prec k$ are both links, then $i \prec k$ is a relation but \emph{not} a link (since $j$ is intermediate), so there is no triangle $i$-$j$-$k$ in the Hasse graph.

Second, the independence number (largest set of mutually non-adjacent vertices in the Hasse diagram) scales as $\alpha \sim N^{0.83}$. Since the Hasse diagram has $\sim N \ln N$ edges and $N$ vertices, the independence number is sub-linear but large, reflecting the sparse, bipartite-like structure of the link graph.

\section{Discussion}

\textbf{1. The Hasse diagram as geometric probe.} Most causal set observables use the transitive closure. We have shown that the Hasse diagram---the minimal, irreducible causal structure---carries complementary geometric information. Its algebraic connectivity sharply distinguishes manifold-like causets from random DAGs, is sensitive to the BD phase transition, and encodes the embedding dimension.

\textbf{2. High connectivity and locality.} The Fiedler value of Minkowski-embedded causets saturates at $\lambda_2 \approx 1.4$--$1.6$ for $N > 100$ (scaling as $N^{0.054}$), much slower than the small-$N$ apparent scaling of $N^{0.32}$. Despite this saturation, the Fiedler value remains substantially higher than that of random DAGs at all $N$ tested. The Hasse diameter reaches 6 at $N = 3000$ and grows as $N^{0.11}$---\emph{not} $\Theta(\sqrt{N})$ as might be expected from the Vershik-Kerov chain scaling. Physically, the high but bounded algebraic connectivity reflects the locality of Minkowski spacetime: every spacetime region is connected to every other by short chains of links, with no bottlenecks. Random DAGs lack this spatial locality, producing fragmented Hasse diagrams with low connectivity. The triangle-free property (girth $\geq 4$) is a direct consequence of transitivity in the partial order.

\textbf{2b. Spectral embedding.} A spectral embedding using the 19 smallest eigenvectors of the Hasse Laplacian achieves $R^2 = 0.83$--$0.91$ for reconstructing embedding coordinates, confirming that the Hasse Laplacian encodes the geometric layout of the causal set.

\textbf{3. The transition restructures links, not relations.} The BD transition changes the Hasse diagram topology while preserving the total number of causal relations. This complements the interval entropy perspective (Paper A), which focuses on the distribution of interval sizes. Together, interval entropy and link fraction provide a complete characterization: the transition concentrates intervals onto direct links (high $\ell$, low $H$).

\textbf{4. Limitation at high dimension.} The Fiedler value collapses to zero for $d \geq 4$ because the Hasse diagram becomes disconnected. This is a fundamental limitation of Hasse-based observables in physically relevant dimensions ($d = 4$). Alternative graph-theoretic probes (component structure, local clustering) may be needed.

\section{Conclusion}

The Hasse diagram of causal sets carries rich geometric information that complements transitive-closure-based observables. The algebraic connectivity is $6$--$80\times$ larger for Minkowski-embedded causets than for random DAGs at small $N$, scaling as $\lambda_2 \sim N^{0.32}$ for $N = 30$--$70$, but \emph{saturating} at $\lambda_2 \approx 1.4$--$1.6$ for $N > 100$ (large-$N$ scaling $\sim\!N^{0.054}$). The Hasse diameter reaches 6 at $N = 3000$, growing as $N^{0.11}$. The Fiedler value jumps by 74\% at the BD phase transition and encodes the embedding dimension via monotonic decrease with $d$. A spectral embedding with 19 eigenvectors achieves $R^2 = 0.83$--$0.91$. The Hasse diagram is triangle-free (girth $\geq 4$) with independence number $\alpha \sim N^{0.83}$. The link fraction provides a perfectly monotonic order parameter for the BD transition, capturing the restructuring of irreducible causal relations. These results establish spectral graph theory on the Hasse diagram as a new tool for causal set quantum gravity, with the caveat that the Fiedler value saturates at large $N$ rather than growing without bound.

\begin{acknowledgments}
This work was conducted with assistance from Claude (Anthropic).
\end{acknowledgments}

\begin{thebibliography}{99}

\bibitem{bombelli1987}
L.~Bombelli, J.~Lee, D.~Meyer, and R.~D.~Sorkin,
``Space-time as a causal set,''
Phys.\ Rev.\ Lett.\ \textbf{59}, 521 (1987).

\bibitem{surya2019}
S.~Surya,
``The causal set approach to quantum gravity,''
Living Rev.\ Rel.\ \textbf{22}, 5 (2019),
\href{https://arxiv.org/abs/1903.11544}{arXiv:1903.11544}.

\bibitem{myrheim1978}
J.~Myrheim,
``Statistical geometry,''
CERN preprint TH-2538 (1978).

\bibitem{meyer1988}
D.~A.~Meyer,
``The dimension of causal sets,''
Ph.D.\ thesis, MIT (1988).

\bibitem{benincasa2010}
D.~M.~T.~Benincasa and F.~Dowker,
``The scalar curvature of a causal set,''
Phys.\ Rev.\ Lett.\ \textbf{104}, 181301 (2010),
\href{https://arxiv.org/abs/1001.2725}{arXiv:1001.2725}.

\bibitem{fiedler1973}
M.~Fiedler,
``Algebraic connectivity of graphs,''
Czech.\ Math.\ J.\ \textbf{23}, 298 (1973).

\bibitem{chung1997}
F.~R.~K.~Chung,
\emph{Spectral Graph Theory} (AMS, 1997).

\bibitem{glaser2018}
L.~Glaser, D.~O'Connor, and S.~Surya,
``Finite size scaling in 2d causal set quantum gravity,''
Class.\ Quant.\ Grav.\ \textbf{35}, 045006 (2018),
\href{https://arxiv.org/abs/1706.06432}{arXiv:1706.06432}.

\bibitem{brightwell1991}
G.~Brightwell and R.~Gregory,
``Structure of random discrete spacetime,''
Phys.\ Rev.\ Lett.\ \textbf{66}, 260 (1991).

\bibitem{loftus_interval}
M.~Loftus,
``Interval entropy as an order parameter for the phase transition in causal set quantum gravity in two and four dimensions,''
(companion paper).

\bibitem{loomis2018}
S.~P.~Loomis and S.~Carlip,
``Suppression of non-manifold-like sets in the causal set path integral,''
Class.\ Quant.\ Grav.\ \textbf{35}, 024002 (2018),
\href{https://arxiv.org/abs/1709.00064}{arXiv:1709.00064}.

\end{thebibliography}

\end{document}
